\section{Teori Relativitas Umum}
\label{sec:relativitas-umum}

\noindent Teori relativitas umum yang dirumuskan oleh Einstein mengubah pemahaman manusia mengenai gravitasi. 
Dalam pandangan Newtonian, gravitasi dipahami sebagai gaya tarik-menarik antar dua benda yang bermassa.
Einstein mengubah pandangan ini dengan memandang gravitasi bukan sebagai gaya, melainkan sebagai 
manifestasi geometris dari kelengkungan ruang-waktu yang disebabkan oleh distribusi materi dan energi 
\parencite{Einstein1916}. Dengan demikian, terdapat hubungan timbal balik materi dan energi menentukan 
bagaimana ruang-waktu dapat melengkung, dan sebaliknya, lengkungan ruang-waktu menentukan bagaimana materi 
dapat bergerak melaluinya.

Dalam relativitas umum, geometri ruang-waktu sepenuhnya ditentukan oleh tensor metrik 
($g_{\mu\nu}$). Tensor metrik merupakan besaran fundamental yang mendefinisikan bagaimana 
jarak dan selang waktu diukur pada setiap titik di ruang-waktu. Dari tensor metrik ini 
diturunkan simbol Christoffel ($\Gamma^\lambda_{\mu\nu}$), yaitu besaran yang mendefinisikan 
aturan untuk membandingkan vektor di dua titik berbeda di ruang-waktu lengkung, karena di 
ruang-waktu yang melengkung, vektor di satu titik tidak dapat langsung dibandingkan dengan 
vektor di titik lain tanpa aturan transportasi yang konsisten. Simbol Christoffel inilah 
yang muncul dalam persamaan geodesik, yaitu persamaan yang mendeskripsikan lintasan alami 
sebuah partikel akibat kelengkungan ruang-waktu, sehingga menjadi penghubung matematis 
antara geometri dan gerak materi yang telah disinggung sebelumnya 
\parencite{Wald1984, Carroll2004}.

Selanjutnya, dari simbol Christoffel, kelengkungan ruang-waktu dapat dikuantifikasi secara 
matematis melalui tensor Riemann \parencite{Wald1984}. Tensor Riemann ($R^\rho_{\sigma\mu\nu}$) memuat seluruh 
informasi kelengkungan secara lengkap di setiap titik dan dari segala arah, namun jumlah 
komponen independennya sangat banyak sehingga tidak praktis untuk digunakan langsung dalam 
persamaan medan. Informasi kelengkungan tersebut dapat diringkas melalui operasi yang disebut kontraksi. 
Dengan operasi ini, informasi kelengkungan yang semula 
mencakup seluruh arah secara rinci dirangkum menjadi informasi yang lebih umum, dengan 
konsekuensi bahwa sebagian detail per arah tidak lagi termuat secara terpisah. 
Kontraksi pertama pada tensor Riemann menghasilkan tensor Ricci 
($R_{\mu\nu}$), yang merangkum informasi kelengkungan per arah di suatu titik. Kontraksi 
lebih lanjut pada tensor Ricci menghasilkan skalar Ricci ($R$), yaitu besaran tunggal tanpa 
indeks yang merangkum kelengkungan keseluruhan di suatu titik tanpa membedakan arah 
\parencite{Carroll2004, Misner1973}.

Hubungan antara geometri kelengkungan ruang-waktu dan distribusi materi-energi
di dalamnya dinyatakan melalui persamaan medan Einstein. Dengan menggunakan
\textit{natural unit} $G = c = 1$, persamaan medan Einstein berbentuk
\begin{equation}
    G_{\mu\nu} + \Lambda g_{\mu\nu} = 8\pi T_{\mu\nu}
    \label{eq:efe}
\end{equation}
\noindent dengan $G_{\mu\nu}$ merupakan tensor Einstein yang merepresentasikan geometri 
ruang-waktu, $\Lambda$ adalah konstanta Kosmologi, dan $T_{\mu\nu}$ adalah tensor 
energi-momentum yang merepresentasikan distribusi materi dan energi \parencite{Carroll2004}. 
Tensor Einstein sendiri dibangun dari tensor Ricci dan skalar Ricci yang telah 
diperkenalkan sebelumnya, yaitu
\begin{equation}
    G_{\mu\nu} = R_{\mu\nu} - \frac{1}{2} R g_{\mu\nu}.
    \label{eq:einstein-tensor}
\end{equation}

Bentuk persamaan ini memperlihatkan secara matematis hubungan timbal balik yang telah 
disinggung di awal: ruas kiri pada Persamaan \eqref{eq:efe} merepresentasikan geometri 
ruang-waktu, sedangkan ruas kanan merepresentasikan kandungan materi dan energi 
\parencite{Schutz2009}. Dengan demikian, distribusi materi dan energi menentukan 
bagaimana ruang-waktu melengkung, dan sebaliknya kelengkungan tersebut menentukan 
bagaimana materi bergerak melaluinya, sebagaimana telah dijelaskan melalui peran simbol 
Christoffel dalam persamaan geodesik. Penurunan lengkap persamaan medan Einstein dapat 
dilihat pada Lampiran~\ref{app:penurunan-efe}.
